\documentclass[10pt,a4paper]{amsart}
\usepackage[margin=19mm]{geometry}
\usepackage{amsmath,amssymb,mathtools}
\usepackage[scr=rsfs,cal=euler]{mathalfa}
\usepackage{xfrac}
\usepackage[unicode,hidelinks,bookmarks=false]{hyperref}
\newcommand{\ie}{i.e.\ }
\newcommand{\cf}{cf.\ }
\newcommand{\resp}{respectively\ }
\newcommand{\Subsection}{\subsection}
\providecommand{\nobreakdash}{\nobreak\hbox{-}}
\setlength{\emergencystretch}{2em}
\multlinegap0pt
\usepackage[shortlabels]{enumitem}
\usepackage{amssymb}
\newtheorem{prop}{Proposition}
\newtheorem{thm}{Theorem}
\DeclareMathOperator{\C}{\mathsf C}
\DeclareMathOperator{\D}{\mathsf D}
\DeclareMathOperator{\F}{\mathsf F}
\newcommand{\FMod}{{\rm FMod}}
\newcommand{\Hom}{{\rm Hom}}
\newcommand{\Ind}{{\rm Ind}}
\newcommand{\Mod}{{\rm Mod}}
\newcommand{\N}{\mathbb{N}}
\DeclareMathOperator{\R}{\mathsf R}
\newcommand{\id}{{\rm id}}
\title{Separable functors and firm modules}
\author{Patrik Lundstr\"om}
\date{Source: arXiv:2602.13417v3 (CC BY 4.0)}
\begin{document}
\maketitle

\noindent\emph{Source and licence.} Separable functors and firm modules. Version arXiv:2602.13417v3 (CC BY 4.0), distributed by arXiv under the Creative Commons Attribution 4.0 International licence, http://creativecommons.org/licenses/by/4.0/, as recorded in the arXiv licence metadata for this exact version and shown on the abstract page https://arxiv.org/abs/2602.13417v3. Full licence notice: \texttt{licenses/arxiv-2602.13417-CC-BY-4.0.txt}.\par\smallskip

\noindent\textit{Editorial scope.} Counted unit: the article's thm thm (source0.tex lines 1789--1797). The article's complete original proof of it is delivered with it (source0.tex lines 1799--1991, one \texttt{proof} environment of 193 lines). No mathematics is written, repaired or re-derived: the statement and the proof are the article's own lines, in the article's own notation, and no retained line is substituted or re-flowed.  Retained uncounted as the source-local context the proof needs: lines 435--453.  Nothing was omitted from any retained span, so the package carries no omission record.  No retained line contains a citation, so the wrapper carries no bibliography and no bibliography entry.  No retained line contains a figure environment, an image-inclusion command or drawing code, so no image is delivered and the package declares no asset dependency.  Editorial formatting only, in addition: an AMS \texttt{amsart} preamble that restates the article's own declarations and macros used by the retained lines, namely the environments \texttt{prop}, \texttt{thm} on separate counters, together with the standard packages those lines need.  The retained environments print in this document as Proposition 1, Theorem 1 in order of appearance, whereas the article prints the same items with the numbering of its own full document.  The complete native source of the article as distributed in the arXiv e-print for this exact version is bundled unchanged as \texttt{sources/194635/source0.tex}.
\begin{prop}\label{prop:newsep}
Let $\F : \C \to \D$ be a functor. 
Then $\F$ is 
separable if and only if for all 
$M,N \in \C_0$, there is a map $\R_{M,N}$ 
satisfying (SF0), (SF1), and 
\begin{itemize}[widest=(SF2)]

\item[{\rm (SF3)}] 
$\R_{M,P}(gf) = 
\R_{N,P}(g) \R_{M,N}(f)$ for $M,N,P \in \C_0$,
$f : \F(M)  \to  \F(N)$, 
\linebreak
$g :  \F(N)  \to  \F(P)$, 
whenever
$f \in \F( \Hom_{\C}(M,N) )$ or 
$g \in \F( \Hom_{\C}(N,P))$.
\end{itemize}
\end{prop}

\begin{thm}
Let $B$ be a firm ring and suppose that 
$f : B \to A$ is a firm 
ring homomorphism. Then the functor
$\Ind_f : {}_B \FMod \to {}_A \FMod$ is 
separable if and only if 
$f$ is a split monomorphism in the category 
of $B$-bimodules. 
\end{thm}

\begin{proof}
First we show the ``only if'' statement.
Suppose that $\Ind_f$ is separable.
We wish to show that $f$ is a split monomorphism
in the category of $B$-bimodules.
Set $\sigma_{A,B} := \mu_{A,B}^{-1}$ 
and $\gamma := (\mu_{A/B} \otimes \id_B) \circ 
(\id_A \otimes \sigma_{A,B})$. 
Then $\gamma$ is 
a morphism in ${}_A \Mod_B$. 
Now we show that
\begin{equation}\label{eq:ABgammaAB}
\mu_{A,B} \circ \gamma \circ 
\Ind_f(f) = 
\mu_{A,B} \circ 
\Ind_f(\id_B). 
\end{equation}
Suppose that 
$a \in A$ and $b \in B$.
Take $n \in \N$, $a_i \in A$ and $b_i \in B$
such that $\sigma_{A,B}( f(b) ) = 
\sum_{i=1}^n a_i \otimes b_i$.
Then $\sum_{i=1}^n a_i f(b_i) = 
\mu_{A,B}(\sigma_{A,B}(f(b))) = f(b)$.
Thus,
\[
\begin{array}{rcl}
(\mu_{A,B} \circ \gamma \circ \Ind_f(f))
(a \otimes b) 
&=& 
\sum_{i=1}^n 
(\mu_{A,B} \circ (\mu_{A/B} \otimes \id_B))
(a \otimes a_i \otimes b_i) \\ 
&=&
\sum_{i=1}^n \mu_{A,B}( a a_i \otimes b_i ) 
= \sum_{i=1}^n a a_i f(b_i) \\
&=& 
a f(b) = 
(\mu_{A,B} \circ \id_{A \otimes_B B})(a \otimes b).
\end{array}
\]
This establishes \eqref{eq:ABgammaAB}.
Since $\mu_{A,B}$ is bijective, we get 
$\gamma \circ \Ind_f(f) = 
\Ind_f(\id_B)$.
By (SF3), 
$\R_{A,B}^{\Ind_f}(\gamma) \circ f = \id_B$.
Thus, $f$ splits as an additive map.
Since $\gamma$ is left $A$-linear,
$\R_{A,B}^{\Ind_f}(\gamma)$ is 
left $B$-linear. Now we show that 
$\R_{A,B}^{\Ind_f}(\gamma)$
is right $B$-linear.
To this end, take $b \in B$ and let
$m_b^B : B \to B$ and $m_b^A : A \to A$
denote the left $B$-linear maps given by right multiplication by $b$ and $f(b)$, respectively.
Since $\gamma$ is right $B$-linear, we have
$\gamma \circ \Ind_f(m_b^A)=\Ind_f(m_b^B)\circ \gamma$.
By (SF3), $\R_{A,B}^{\Ind_f}(\gamma)\circ m_b^A
= m_b^B\circ \R_{A,B}^{\Ind_f}(\gamma)$.
Therefore, $\R_{A,B}^{\Ind_f}(\gamma)$ is right $B$-linear.

Now we show the ``if'' statement.
Suppose that $f$ is a split monomorphism in the 
category of $B$-bimodules. Then there is a 
$B$-bimodule map
$s : A \to B$ such that $s \circ f = \id_B$.
Suppose that $M$ and $N$ are firm left $B$-modules.
Define 
$\R_{M,N}^{\Ind_f} : 
\Hom_A(\Ind_f(M),\Ind_f(N)) \to 
\Hom_B(M,N)$
in the following way. 
Take 
$g \in \Hom_A(\Ind_f(M) , \Ind_f(N))$.
Put
\[
\R_{M,N}^{\Ind_f}(g) :=
\mu_{B,N} \circ (s \otimes \id_N) 
\circ g \circ 
(f \otimes \id_M) \circ \mu_{B,M}^{-1}.
\]
Then $\R_{M,N}^{\Ind_f}(g)$
is a morphism in ${}_B \Mod$.
Now we show (SF1).
Take $m \in M$,
$k \in \N$, $b_i \in B$ and $m_i \in M$
with $\sum_{i=1}^k b_i m_i = m$. 
Let $g = \Ind_f(g')$ for some 
$g' \in \Hom_B(M,N)$. Then 
\[
\begin{array}{rcl}
\R_{M,N}^{\Ind_f}(g)(m) 
&=& 
\sum_{i=1}^k (\mu_{B,N} \circ
(s \otimes \id_N) \circ g)
(f(b_i) \otimes m_i) \\
&=&
\sum_{i=1}^k
(\mu_{B,N} \circ (s \otimes \id_N)
\circ (\id_A \otimes g')) 
(f(b_i) \otimes m_i) \\ 
&=& 
\sum_{i=1}^k
(\mu_{B,N} \circ (s \otimes \id_N))
(f(b_i) \otimes g'(m_i)) \\ 
&=& 
\sum_{i=1}^k
\mu_{B,N} (b_i \otimes g'(m_i)) \\
&=& \sum_{i=1}^k b_i g'(m_i) 
= g' \left( \sum_{i=1}^k b_i m_i \right)
= g'(m).
\end{array}
\]
Now we show (SF3).
Suppose that $M,N,P$ are modules in ${}_B \FMod$. 
Take $g \in \Hom_A(A \otimes_B N,
A \otimes_B P)$ and 
$h \in \Hom_A(A \otimes_B M,
A \otimes_B N)$. 
We consider two cases.

Case 1: $g \in \Ind_f(\Hom_B(N,P))$
that is 
$g = \id_A \otimes g'$ for 
some $g' \in \Hom_B(N,P)$. Then
\[
(s \otimes \id_P) \circ g \circ 
(f \otimes \id_N) \circ 
(s \otimes \id_N) = 
(\id_B \otimes g') \circ 
(s \otimes \id_N) = 
(s \otimes \id_P) \circ g.
\]
Therefore,
{\small \[
\begin{array}{cl}
 & \R_{N,P}^{\Ind_f}(g) \circ
\R_{M,N}^{\Ind_f}(h) \\ 
= & \mu_{B,P} \circ (s \otimes \id_P)
\circ g \circ (f \otimes \id_N)
\circ \mu_{B,N}^{-1} \, \circ 
 \mu_{B,N} \circ (s \otimes \id_N)
\circ h \circ (f \otimes \id_M) 
\circ \mu_{B,M}^{-1} \\
=& \mu_{B,P} \circ (s \otimes \id_P)
\circ g \circ (f \otimes \id_N) \, \circ   
(s \otimes \id_N) \circ h \circ 
(f \otimes \id_M) \circ \mu_{B,M}^{-1} \\
=& \mu_{B,P} \circ (s \otimes \id_P) 
\circ g \circ h \circ 
(f \otimes \id_M) \circ \mu_{B,M}^{-1} 
= \R_{M,P}^{\Ind_f}(g \circ h).
\end{array}
\]}

Case 2: $h \in \Ind_f(\Hom_B(M,N))$
that is $h = \id_A \otimes h'$ for 
some $h' \in \Hom_B(M,N)$. Then
\[
(f \otimes \id_N) \circ
(s \otimes \id_N) \circ h \circ 
(f \otimes \id_M)  = 
(\id_A \otimes h') \circ 
(s \otimes \id_M) = 
h \circ (f \otimes \id_M).
\]
Therefore,
{\small \[
\begin{array}{cl}
 & \R_{N,P}^{\Ind_f}(g) \circ
\R_{M,N}^{\Ind_f}(h) \\
=& \mu_{B,P} \circ (s \otimes \id_P)
\circ g \circ (f \otimes \id_N)
\circ \mu_{B,N}^{-1} \, \circ 
\mu_{B,N} \circ (s \otimes \id_N)
\circ h \circ (f \otimes \id_M) 
\circ \mu_{B,M}^{-1} \\
=& \mu_{B,P}  \circ  (s \otimes \id_P)
\circ g \circ 
(f \otimes \id_N) \, \circ 
(s \otimes \id_N)  \circ h \circ
(f \otimes \id_M)  \circ \mu_{B,M}^{-1} \\
=& \mu_{B,P} \circ (s \otimes \id_P) 
\circ g \circ h \circ 
(f \otimes \id_M) \circ \mu_{B,M}^{-1} 
= \R_{M,P}^{\Ind_f}(g \circ h).
\end{array}
\]}
This establishes (SF3).
By Proposition \ref{prop:newsep},
$\Ind_f$ is separable.
\end{proof}

\end{document}
